\section{Templates}

% Motivating example
\begin{frame}[fragile]{Motivation: \texttt{MaybeInt}}
  Suppose we want a type that either holds a value, or nothing (like a nullable int):
  \begin{minted}{cpp}
    struct MaybeInt {
        bool mHasValue; int mValue;

        MaybeInt() : mHasValue(false), mValue{} {}
        MaybeInt(int v) : mHasValue(true), mValue(v) {}

        bool has() const { return mHasValue; }
        int get() const {
            if (!mHasValue) throw std::runtime_error("empty");
            return mValue;
        }
    };
  \end{minted}
  What if we want \texttt{MaybeDouble}? \texttt{MaybeString}? Same logic, different type.
\end{frame}

% Class templates
\begin{frame}[fragile]{\texttt{Maybe} - Templated Solution}
  \begin{sqt:result}
    % Recall \texttt{MaybeInt} from earlier. \\
    We can generate the exact same code by writing a \emph{template class}.
  \end{sqt:result}

  \begin{minted}{cpp}
    /* 4.02 */
    template <typename T>
    struct Maybe {
        bool mHasValue; T mValue;
        /* ... */
    };
  \end{minted}

  Verify this with \verb|cppinsights.io|. \\
  % \textit{Bonus question: what's the difference between \texttt{struct} and \texttt{class}?}
\end{frame}

\begin{frame}[fragile]{Template Instantiation}
  \begin{columns}[T]
    \begin{column}{0.4\textwidth}
      \begin{minted}{cpp}
  /* 4.02 - CPPINSIGHTS OUTPUT */

#ifdef INSIGHTS_USE_TEMPLATE
template<> struct Maybe<int> {
  bool mHasValue;
  int mValue;
  /* ... */
  inline const int & get() // ...
};
#endif

#ifdef INSIGHTS_USE_TEMPLATE
template<> struct Maybe<double> {
  bool mHasValue;
  double mValue;
  /* ... */
  inline const double & get() //...
};
#endif
\end{minted}
    \end{column}
    \begin{column}{0.6\textwidth}
      \vspace{1cm}
      For each type (\texttt{int, std::string, double}) that is substituted for
      the template type parameter \texttt{T}, 
      an \emph{instantiation} is generated.

      Based on the code, the compiler determines which instantiations are necessary,
      and will only generate those which are used.
    \end{column}
  \end{columns}
\end{frame}

% Template syntax
\begin{frame}[fragile]{C++ Templates}
  \begin{minted}{cpp}
// declare Widget as a template class
template <typename T, typename U, int V> struct Widget{/* */};
// declare void foo(T,U) as a template function
template <typename T, typename U> void foo(T lorem, U ipsum) {/* */}

int main() {
    auto s = std::string{"hi"};
    auto val = 42;
    foo(s, val); // instantiation <T = std::string, U = int>
    auto obj = Widget<int, decltype(s), 20>{}; // instantiation <T = int, U = decltype(s), V = 20>
    auto obj2 = Widget<int, int, val>{}; // COMPILE ERROR: val is not a compile-time value
}
  \end{minted}

  \begin{sqt:result}
    The programmer can \emph{use} the template entity (function / class).
    At compile-time, the compiler \emph{stamps} out every required instantiation, \\
    replacing each template parameter with the required types / compile-time values.
  \end{sqt:result}
\end{frame}

\begin{frame}[fragile]{What can be templated?}
  Templates can be used to define:
  \begin{itemize}
    \item \textbf{Function templates}: generate function definitions, with types substituted
    \item \textbf{Class templates}: generate class definitions, with types substituted.
    \item \textbf{Variable templates}: generate variable definitions, with types substituted
  \end{itemize}

  Let's take a look at some more examples.
\end{frame}

% Function templates
\begin{frame}[fragile]{Example: \texttt{max()} without function templates}
  Suppose we are writing our own custom \code{cmath} library:
  \begin{minted}{cpp}
    int max(int x, int y) {
        return (x > y) ? x : y;
    }
    double max(double x, double y) {
        return (x > y) ? x : y;
    }
    char max(char x, char y) {
        return (x > y) ? x : y;
    }

    int main() {
        std::cout << max(1, 2) << '\n';     // overload resolution: max(int, int)
        std::cout << max(3.5, 7.5) << '\n'; // overload resolution: max(double, double)
        std::cout << max('a', 'b') << '\n'; // overload resolution: max(char, char)
        return 0;
    }
  \end{minted}

  Can we do better?
\end{frame}

\begin{frame}[fragile]{Example: \texttt{max()} with function templates}
  \begin{minted}{cpp}
    #include <iostream>

    template <typename T>
    T max(T x, T y) {
        return (x > y) ? x : y;
    }

    int main() {
        std::cout << max(1, 2) << '\n';     // deduces to max<int>(int)
        std::cout << max(3.5, 7.5) << '\n'; // deduces to max<double>(double)
        std::cout << max('a', 'b') << '\n'; // deduces to max<char>(char)
        return 0;
    }
  \end{minted}

  What about \texttt{max(1, 2.1)}?
\end{frame}

\begin{frame}[fragile]{Example: Multiple type parameters}
  \begin{minted}{cpp}
    #include <iostream>

    template <typename T, typename U>
    auto max(T x, U y) {
        return (x > y) ? x : y;
    }

    int main() {
        std::cout << max(1, 2.1) << '\n'; // 2.1
        return 0;
    }
  \end{minted}

  Using \texttt{auto}, the compiler deduces the return type.
\end{frame}

\begin{frame}[fragile]{Where have we seen function templates?}
  \texttt{std::sort} from \texttt{<algorithm>} works on any container of any type:
  \begin{minted}{cpp}
    #include <algorithm>
    #include <vector>

    auto ints = std::vector<int>{3, 1, 2};
    auto doubles = std::vector<double>{3.5, 1.1, 2.2};
    auto strings = std::vector<std::string>{"banana", "apple", "cherry"};

    std::sort(ints.begin(), ints.end());
    std::sort(doubles.begin(), doubles.end());
    std::sort(strings.begin(), strings.end());
  \end{minted}
  Under the hood, \texttt{std::sort} is a function template; same code,
  ``stamped out" once for each type it is called with.
\end{frame}

\begin{frame}[fragile]{Template instantiation examples}
  \begin{minted}{cpp}
    /* 4.03 */
    // an array of 8 integers
    struct ArrayInt8 {
        using T = int;
        static constexpr std::size_t N = 8;
        T mData[N];
        // set array element to a value by index
        void set(std::size_t index, T val) {/* */}
        // get const reference to array element by index
        const T& at(std::size_t index) const {/* */}
    };

    // returns two values: {maximum in nums, average of nums}
    Pair<int, double> getMaxAndAverage(const ArrayInt8& nums) {
        /* ... */
        return {currMax, currSum / 8};
    }
  \end{minted}

  Copy the full code example into cppinsights to see all generated instantiations.
\end{frame}

\begin{frame}[fragile]{Templating the array class}
  \begin{minted}{cpp}
    /* 4.04 */
    template <typename T, std::size_t N>
    struct Array {
        T mData[N];
        // set array element to a value by index
        void set(std::size_t index, T val) {/* */}
        // get const reference to array element by index
        const T& at(std::size_t index) const {/* */}
    };

    // returns two values: {maximum in nums, average of nums}
    Pair<int, double> getMaxAndAverage(const Array<int, 8>& nums) {
        /* ... */
        return {currMax, currSum / 8};
    }
  \end{minted}

  Copy the full code example into cppinsights to see all generated instantiations.
\end{frame}

\begin{frame}[fragile]{Templating the function}
  \begin{minted}{cpp}
    /* 4.05 */
    template <typename T, std::size_t N>
    // requires: T - convertible to double
    Pair<T, double> getMaxAndAverage(const Array<T, N>& nums) {
        auto currMax = nums.at(0);
        auto currSum = 0.0;
        for (std::size_t i = 0; i < N; ++i) {
            currMax = currMax > nums.at(i) ? currMax : nums.at(i);
            currSum += nums.at(i);
        }
        return {currMax, currSum / N};
    }
  \end{minted}

  Copy the full code example into cppinsights to see all generated instantiations.
\end{frame}

\begin{frame}[fragile]{Instantiation Failures}
  \begin{minted}{cpp}
    /* 4.06 */
    template <typename Container, typename Value>
    Value pop(Container& c) {
        c.pop_front();
        /* ... */
    }

    /* compiler generated instantiation */ /*Container = double, Value = int*/
    int pop<double, int>(double& c) {
        c.pop_front(); // COMPILE ERROR!
        /* ... */
    }

    int main() {
        auto x = double{};
        pop(x, 1);
    }
  \end{minted}
\end{frame}

\begin{frame}[fragile]{Exercise: Queue}
  You are given an implementation of a queue of \code{int}.\\
  Make it generic.

  \begin{minted}{cpp}
    /* 4.08 */
    // given a queue of ints
    class Queue {/* ... */};

    // implement a queue of "any" type
    template <typename T>
    class Queue {/* ... */};
  \end{minted}
\end{frame}